%=============================================================================!
% 11.5  RESONANCE                                                             !
%=============================================================================!
\section{Resonance}
\label{sec:cs-resonance}

RESPA applies the slow forces as sudden kicks, once per outer step, to
atoms that are vibrating under the fast forces in between. A swing pushed
in step with its motion goes higher, and without drag higher and higher
(\cref{sec:os-driven}); kicks that come in step with a vibration, once in
every half swing or every swing, can pump it in the same way. This section
finds the outer steps at which that happens for the simplest system that
shows it, and compares the result with water.

\subsection*{Two springs}

One coordinate $q$, of unit mass, is pulled back by a fast spring and a
slow one, with angular frequencies $\omega_{\mathrm f}$ and
$\omega_{\mathrm s}$, so its force is $-(\omega_{\mathrm f}^2 +
\omega_{\mathrm s}^2)\,q$. RESPA follows the fast spring with many inner
steps, closely enough to take its motion as exact: by the solution of
\cref{sec:ne-spring}, over a time $\dt$ it carries the state $(q, v)$ to
\begin{equation*}
   q' = q\cos\theta + \frac{v}{\omega_{\mathrm f}}\sin\theta , \qquad
   v' = -\omega_{\mathrm f}q\sin\theta + v\cos\theta , \qquad
   \theta = \omega_{\mathrm f}\dt ,
\end{equation*}
and half a kick by the slow spring changes $v$ by $-\frac12\dt\,
\omega_{\mathrm s}^2q$ and leaves $q$. Both are linear, so as matrices
acting on the column $(q, v)$,
\begin{equation*}
   \vb{R} = \begin{pmatrix} \cos\theta & \sin\theta/\omega_{\mathrm f}\\
   -\omega_{\mathrm f}\sin\theta & \cos\theta \end{pmatrix} , \qquad
   \vb{K} = \begin{pmatrix} 1 & 0\\ -\kappa & 1 \end{pmatrix} , \qquad
   \kappa = \tfrac12\dt\,\omega_{\mathrm s}^2 ,
\end{equation*}
and the outer step is $\vb{M} = \vb{K}\vb{R}\vb{K}$, the half kick
$\vb{K}$ acting first. Multiplying, with $C = \cos\theta$ and $S =
\sin\theta$,
\begin{align*}
   \vb{R}\vb{K} &= \begin{pmatrix} C - \kappa S/\omega_{\mathrm f} &
      S/\omega_{\mathrm f}\\ -\omega_{\mathrm f}S - \kappa C & C
      \end{pmatrix} ,\\
   \vb{K}\vb{R}\vb{K} &= \begin{pmatrix} C - \kappa S/\omega_{\mathrm f} &
      S/\omega_{\mathrm f}\\ -\omega_{\mathrm f}S - 2\kappa C
      + \kappa^2S/\omega_{\mathrm f} & C - \kappa S/\omega_{\mathrm f}
      \end{pmatrix} ,
\end{align*}
so the trace, the sum of the diagonal entries (\cref{sec:in-reversal}),
is $2C - 2\kappa S/\omega_{\mathrm f}$, which with $2\kappa = \dt\,
\omega_{\mathrm s}^2$ is
\begin{equation}
   \operatorname{tr}\vb{M} = 2\cos\theta
   - \frac{\omega_{\mathrm s}^2\dt}{\omega_{\mathrm f}}\sin\theta .
   \label{eq:cs-trace}
\end{equation}
The fast motion $\vb{R}$, which turns the point $(q, v/\omega_{\mathrm f})$
through the angle $\theta$, and the kick $\vb{K}$ each keep area in the
phase plane (determinants $\cos^2\theta + \sin^2\theta = 1$ and $1$), so
their product does too, and $\det\vb{M} = 1$ (\cref{sec:ha-symplectic}).
Whether repeated outer steps stay bounded depends on the trace alone.

\begin{toolbox}[Powers of a $2\times2$ matrix]
Any $2\times2$ matrix $\vb{M}$, with $M_{11}$, $M_{12}$ in its first row
and $M_{21}$, $M_{22}$ in its second, satisfies
\begin{equation*}
   \vb{M}^2 = (\operatorname{tr}\vb{M})\,\vb{M} - (\det\vb{M})\,\vb{I} ,
\end{equation*}
as multiplying out shows: $\vb{M}^2$ has $M_{11}^2 + M_{12}M_{21}$ and
$M_{11}M_{12} + M_{12}M_{22}$ in its first row and $M_{21}M_{11} +
M_{22}M_{21}$ and $M_{21}M_{12} + M_{22}^2$ in its second, while $(M_{11}
+ M_{22})\vb{M}$ has $M_{11}^2 + M_{11}M_{22}$, $M_{11}M_{12} +
M_{12}M_{22}$ and $M_{21}M_{11} + M_{21}M_{22}$, $M_{11}M_{22} +
M_{22}^2$, and subtracting $(M_{11}M_{22} - M_{12}M_{21})\vb{I}$ from the
latter gives the former. With $\det\vb{M} = 1$, multiplying by
$\vb{M}^{n-1}$ gives $\vb{M}^{n+1} = (\operatorname{tr}\vb{M})\vb{M}^n -
\vb{M}^{n-1}$, so each entry $x_n$ of the state after $n$ steps,
$\vb{M}^n$ times the start, obeys
\begin{equation*}
   x_{n+1} = (\operatorname{tr}\vb{M})\,x_n - x_{n-1} .
\end{equation*}
This is the recurrence of \cref{sec:in-shadow}. If
$\lvert\operatorname{tr}\vb{M}\rvert < 2$, then $\operatorname{tr}\vb{M} =
2\cos\phi$ for some angle $\phi$, and $\cos n\phi$ and $\sin n\phi$ solve
it, as shown there: the motion stays bounded. If
$\lvert\operatorname{tr}\vb{M}\rvert > 2$, the quadratic $\Lambda^2 -
(\operatorname{tr}\vb{M})\Lambda + 1 = 0$ has, by the quadratic formula, a
real root $\Lambda$ with $\lvert\Lambda\rvert > 1$, and dividing it by
$\Lambda$ gives $\Lambda + 1/\Lambda = \operatorname{tr}\vb{M}$. Then
$\Lambda^{n+1} = (\Lambda + 1/\Lambda)\Lambda^n - \Lambda^{n-1}$, so $x_n =
\Lambda^n$ solves the recurrence, as does $\Lambda^{-n}$, and a general
start, a mixture of the two, grows by the factor $\lvert\Lambda\rvert$ at
every step. At $\lvert\operatorname{tr}\vb{M}\rvert = 2$ exactly the growth,
if any, is slower, by a fixed amount at each step.
\end{toolbox}

Without the slow spring, $\operatorname{tr}\vb{M} = 2\cos\theta$, which
never leaves $[-2, 2]$: exact inner motion is stable at any outer step.
The slow spring adds $-b\theta\sin\theta$ to the trace, with $b =
(\omega_{\mathrm s}/\omega_{\mathrm f})^2$, since $\omega_{\mathrm s}^2
\dt/\omega_{\mathrm f} = b\,\omega_{\mathrm f}\dt$. Just below
$\theta = \pi$ we write $\theta = \pi - \epsilon$ with $\epsilon$ small and
positive; then $\cos\theta = -\cos\epsilon \approx -1 + \frac12\epsilon^2$
and $\sin\theta = \sin\epsilon \approx \epsilon$ (\cref{sec:mo-taylor}),
and
\begin{align*}
   \tfrac12\operatorname{tr}\vb{M} &\approx -1 + \tfrac12\epsilon^2
   - \tfrac12b\pi\epsilon && \text{with $\theta \approx \pi$ in the last
   term,}\\
   &< -1 \quad\text{when}\quad \tfrac12\epsilon^2 < \tfrac12b\pi\epsilon ,
   \text{ that is, } 0 < \epsilon < b\pi .
\end{align*}
The motion therefore grows for outer steps in a band just below $\theta =
\pi$, half the fast period $\mathcal{T}_{\mathrm f} = 2\pi/\omega_{\mathrm
f}$, of width $b\pi$ in $\theta$, or $\frac12b\,\mathcal{T}_{\mathrm f}$
in $\dt$; just below $\theta = 2\pi$, where $\sin\theta$ is negative, the
same argument gives a band twice as wide below $\mathcal{T}_{\mathrm f}$
(\cref{ex:cs-band}). Near $\dt = \frac12\mathcal{T}_{\mathrm f}$ the kicks
come every half vibration, when the fast motion has carried $q$ to minus
its value, and each one, proportional to $q$, meets the vibration in the
same phase as the last: a resonance. Together the two springs shorten the
period to $\mathcal{T}_{\mathrm f}/\sqrt{1 + b}$, so the band lies just
below $\frac12\mathcal{T}_{\mathrm f}$, around half that period. At
$\dt = \frac12\mathcal{T}_{\mathrm f}$ exactly, $\theta = \pi$, the trace
is $-2$, and $\vb{M}$ has rows $(-1, 0)$ and $(2\kappa, -1)$: $q$ comes
back to minus itself at every step while $v$ gains $2\kappa q$ each time,
so the motion grows steadily rather than by a factor.

\Cref{fig:cs-resonance} measures it for $\omega_{\mathrm s} =
0.3\,\omega_{\mathrm f}$, so $b = 0.09$ and the predicted width below
$\frac12\mathcal{T}_{\mathrm f}$ is $0.045\,\mathcal{T}_{\mathrm f}$. The
trace of \cref{eq:cs-trace} leaves $[-2, 2]$ for $\dt$ from 0.459 to
$0.500\,\mathcal{T}_{\mathrm f}$, a width of 0.041, since the term
$\frac12b\epsilon^2$ dropped with $\theta \approx \pi$ narrows the band by
the factor $1/(1 + b)$, and from 0.919 to $1.000\,\mathcal{T}_{\mathrm
f}$. RESPA run for 300 steps from $q = 1$ at rest reaches a size
$\sqrt{q^2 + (v/\omega_{\mathrm f})^2}$ of $\num{2.9e17}$ at
$0.48\,\mathcal{T}_{\mathrm f}$, grows steadily to 23 at
$0.50\,\mathcal{T}_{\mathrm f}$, and stays near its start, below 2, at
0.45 and $0.52\,\mathcal{T}_{\mathrm f}$.

\begin{figure}[htb]
\centering
\includegraphics{ch11_constraints/resonance}
\caption{Resonance of a fast spring kicked by a slow one, $\omega_{\mathrm
s} = 0.3\,\omega_{\mathrm f}$, stepped by RESPA with 50 inner steps.
(a) Half the trace of the outer step, \cref{eq:cs-trace}, against the
outer step in units of the fast period; the motion grows where it passes
$\pm1$ (dotted), in the shaded bands, and the dashed curve, $\cos\theta$,
is half the trace without the slow spring. (b) The largest size of the
state, $\sqrt{q^2 + (v/\omega_{\mathrm f})^2}$, in 300 outer steps from $q
= 1$ at rest.}
\label{fig:cs-resonance}
\end{figure}

\subsection*{In water}

Water has many fast vibrations, and its slow forces are not springs, but
the same picture explains part of \cref{fig:cs-respa}. The energy drifts
most for outer steps around half the periods of the model's two
stretches, 3.67 and \SI{4.44}{\femto\second}: by $+0.19$, $+1.14$,
$+2.25$ and \SI{+0.33}{\electronvolt\per\pico\second} at 3.5, 3.75, 4 and
\SI{4.25}{\femto\second}, and the run blows up at
\SI{4.5}{\femto\second}. The fluctuation grows too fast well before those
steps. From an outer step of 1 to \SI{2}{\femto\second} it grows by a
factor of 4.8, a little more than the factor of 4 of $\dt^2$, and from 2
to \SI{3}{\femto\second} by 4.9, against 2.25. Ma, Izaguirre and Skeel
found the same in flexible water and proteins, and traced it to
resonances of a kind that the linear model above lacks, with
instabilities when the outer step is a third or possibly a quarter of the
fastest period\cite{Ma2003}: here 2.44 and \SI{1.83}{\femto\second}. The
outer step of RESPA is therefore kept below a third of the fastest
period, and for this water a quarter, whatever the cost of the slow
forces.

\begin{keyidea}
Kicks by the slow forces once per outer step can pump a fast vibration.
For two springs the outer step's matrix has the trace $2\cos\theta -
(\omega_{\mathrm s}^2\dt/\omega_{\mathrm f})\sin\theta$, and the motion
grows where it leaves $[-2, 2]$, in bands just below multiples of half
the fast period. In water the fluctuation climbs faster than $\dt^2$
beyond about a quarter of the fastest period, where Ma, Izaguirre and
Skeel found resonances beyond this linear picture.
\end{keyidea}

\notebook{11}{move the outer step across the resonance bands}

\begin{selfcheck}
For $\omega_{\mathrm s} = 0.1\,\omega_{\mathrm f}$ and a fast period of
\SI{9}{\femto\second}, where is the band of growth below half the period,
and how wide is it?
\end{selfcheck}
